\chapter{Fondamenti teorici e stato dell'arte}
\label{chap:fondamenti_stato_arte}

Questo capitolo raccoglie i fondamenti necessari a comprendere la progettazione di uno Spiking Transformer compatto per dati neuromorfici. L'esposizione segue il cammino dell'informazione, dal sensore alla sua rappresentazione discreta, dalla memoria del singolo neurone alle relazioni tra feature, fino al costo computazionale e di memoria che ogni scelta comporta.

La sezione~\ref{sec:rappresentazione_event_camera} affronta il primo passaggio obbligato, la conversione del flusso asincrono di eventi in tensori regolari, e mostra come questa operazione, apparentemente preliminare, costituisca già una forma di modellazione. La sezione~\ref{sec:dinamica_neuronale_memoria} descrive la memoria intrinseca dei neuroni spiking e i limiti di uno stato che riassume ricorsivamente il passato. Le sezioni~\ref{sec:snn_spiking_transformer} e~\ref{sec:qkformer_riferimento} ricostruiscono l'incontro tra attenzione e computazione spiking fino a QKFormer, l'architettura da cui deriva la baseline compatta, distinguendo la memoria dei neuroni dalle interazioni esplicite tra istanti diversi. La sezione~\ref{sec:strategie_modellazione_temporale} esamina due modi di rendere esplicito l'accesso alla storia, l'attenzione spazio-temporale e il mixing mediante ritardi, mentre la sezione~\ref{sec:efficienza_neuromorfica_vincoli} introduce le grandezze con cui se ne misura il costo.

Non si tratta di una rassegna esaustiva. Ogni contributo discusso è stato selezionato perché orienta una scelta progettuale compiuta nel seguito del lavoro, ed è presentato in funzione di essa.


% ==============================================================================
%  Rappresentazione dell'informazione nelle event camera
% ==============================================================================

\section{Event camera e rappresentazione degli eventi}
\label{sec:rappresentazione_event_camera}

Una fotocamera convenzionale scatta fotografie a intervalli regolari, per esempio trenta al secondo, e ogni fotografia registra la luminosità di tutti i pixel, anche di quelli in cui non è cambiato nulla. Una event camera funziona in modo opposto. Ogni pixel resta in silenzio finché la luce che riceve non cambia abbastanza, e in quel momento segnala il cambiamento con un \emph{evento}. Una scena immobile non produce quasi nulla, mentre una mano che si muove produce una scia di eventi lungo i propri contorni, nell'istante esatto in cui si sposta.

Per elaborare questo segnale con una rete neurale, che lavora su tensori di dimensioni fisse, occorre raccogliere gli eventi in una struttura regolare. Questa sezione descrive prima il segnale prodotto dal sensore e poi le scelte necessarie per trasformarlo in un tensore. Sono scelte tutt'altro che neutre, perché decidono quali informazioni arrivano alla rete e quali vanno perse prima ancora che l'elaborazione cominci.

\subsection{Dal contrasto luminoso al flusso di eventi}
\label{subsec:contrasto_eventi}

Nel modello ideale di un Dynamic Vision Sensor (DVS) ogni pixel lavora indipendentemente dagli altri e osserva il logaritmo dell'intensità luminosa che lo colpisce, $L(x,y,t)=\log I(x,y,t)$~\cite{lichtsteiner2008dvs}. Il pixel conserva un valore di riferimento, cioè il livello di $L$ al momento del suo ultimo evento, e lo confronta di continuo con il valore attuale. Quando la differenza raggiunge una soglia $C>0$, verso l'alto o verso il basso, il pixel emette un evento e aggiorna il riferimento. In formule, l'evento $i$-esimo viene emesso nell'istante $t_i$ in cui
\begin{equation}
\label{eq:soglia_contrasto}
L(x_i,y_i,t_i)-L(x_i,y_i,t_i^{\mathrm{ref}})=p_i\,C,
\end{equation}
dove $t_i^{\mathrm{ref}}$ è l'istante a cui risale il riferimento, cioè l'evento precedente dello stesso pixel o la sua inizializzazione, e la polarità $p_i$ vale $+1$ se la luce è aumentata e $-1$ se è diminuita~\cite{gallego2022event}. Ogni evento è quindi descritto da quattro numeri,
\begin{equation}
\label{eq:evento_quadrupla}
e_i=(x_i,y_i,t_i,p_i),
\end{equation}
la colonna e la riga del pixel, l'istante e la polarità. Un evento con $x_i=64$, $y_i=30$, $t_i=1\,250\,400\,\mu\mathrm{s}$ e $p_i=+1$, per esempio, dice che il pixel della colonna 64 e della riga 30 è diventato più luminoso $1{,}25$ secondi dopo l'inizio della registrazione.

La figura~\ref{fig:pixel_event_generation} mostra il meccanismo su un singolo pixel, attraversato prima da un oggetto chiaro e poi da uno scuro. Finché la luce cresce, ogni incremento pari a $C$ produce un evento positivo, quando cala ogni decremento pari a $C$ ne produce uno negativo. Due aspetti meritano attenzione. La soglia agisce sulla variazione accumulata dall'ultimo evento, non sulla velocità del cambiamento, sicché nel modello ideale una variazione lenta e una rapida della stessa entità producono lo stesso numero di eventi, soltanto distribuiti in tempi diversi. Le oscillazioni più piccole della soglia, invece, non producono nulla.

\begin{figure}[tbp]
    \centering
    \includegraphics[width=\textwidth]{images/cap1/cap1_generazione_eventi.pdf}
    \caption[Generazione degli eventi in un pixel ideale]{Generazione degli eventi in un pixel ideale, su un segnale sintetico con soglia $C=0{,}2$. In alto, il logaritmo dell'intensità (linea continua), il valore di riferimento del pixel (tratteggio) e la fascia di ampiezza $\pm C$ attorno a esso. Ogni volta che la curva esce dalla fascia il pixel emette un evento, positivo (triangolo rosso) o negativo (triangolo blu), e il riferimento si sposta di $C$. In basso, la sequenza di eventi che il pixel trasmette. L'oscillazione centrale, più piccola della soglia, non produce alcun evento.}
    \label{fig:pixel_event_generation}
\end{figure}

Perché proprio il logaritmo? Perché la differenza tra due logaritmi è il logaritmo del rapporto,
\begin{equation}
\label{eq:rapporto_luminanza}
\Delta L=\log\frac{I(x_i,y_i,t_i)}{I(x_i,y_i,t_i^{\mathrm{ref}})},
\end{equation}
e il pixel reagisce quindi alle variazioni relative della luce, non a quelle assolute. Con una soglia $C=0{,}2$, scelta qui a titolo illustrativo, un evento positivo richiede un aumento della luce di circa il $22\%$, poiché $e^{0{,}2}\approx1{,}22$. Vale allo stesso modo per un pixel in piena luce che passa da $1000$ a $1220$ unità e per uno in penombra che passa da $10$ a $12{,}2$. La stessa scena produce così eventi simili in condizioni di illuminazione molto diverse, ed è una delle ragioni dell'ampia gamma dinamica di questi sensori~\cite{gallego2022event}.

La polarità, va precisato, dice soltanto se in quel pixel la luce è aumentata o diminuita, non in quale direzione si muova qualcosa. Si pensi a una barra scura che scorre verso destra su uno sfondo chiaro. I pixel raggiunti dal suo bordo anteriore si scuriscono ed emettono eventi negativi, quelli lasciati dal bordo posteriore tornano chiari ed emettono eventi positivi. Se la barra scorresse verso sinistra, ogni pixel emetterebbe esattamente la stessa sequenza, e la direzione si potrebbe leggere soltanto nell'ordine con cui si attivano pixel vicini. Allo stesso modo, un evento non implica necessariamente un movimento, perché anche una lampada che sfarfalla fa cambiare la luce di molti pixel insieme. Il sensore reale, infine, si allontana dal modello ideale in vari modi. Le soglie positiva e negativa possono differire tra loro e variare da pixel a pixel, e il rumore, la banda limitata del fotorecettore e un breve periodo refrattario dopo ogni evento aggiungono altri scostamenti~\cite{gallego2022event}.

Una registrazione è dunque una lunga sequenza di eventi in ordine di tempo,
\begin{equation}
\label{eq:stream_completo}
\mathcal E=\{e_i\}_{i=1}^{N},\qquad t_i\leq t_{i+1}.
\end{equation}
Il numero $N$ degli eventi non dipende dalla durata ma da quanto accade nella scena. Nei dataset usati in questo lavoro, per esempio, una parola pronunciata in circa un secondo produce in media poco più di diecimila eventi, mentre un'immagine di CIFAR10-DVS fatta muovere davanti al sensore per un tempo simile ne produce circa duecentomila (sezione~\ref{sec:dataset}). Anche all'interno della stessa registrazione possono alternarsi lunghi silenzi e raffiche improvvise, e chi elabora il segnale deve saper gestire gli uni e le altre.

\subsection{Discretizzazione e significato del clock}
\label{subsec:discretizzazione_clock}

Una rete neurale spiking elabora l'ingresso a passi discreti. Per presentarle un flusso di eventi bisogna quindi decidere che cosa entri in ciascun passo, cioè come tagliare il flusso in pezzi. La regola che associa ogni passo a un intervallo di tempo, o più in generale a un gruppo di eventi, prende il nome di \emph{clock} dell'input. Le regole possibili sono diverse, e la figura~\ref{fig:partizioni_flusso} le confronta sullo stesso flusso.

La più semplice divide il tempo in intervalli consecutivi della stessa durata $\Delta t$, detti bin, a partire da un'origine $t_0$. Il bin $k$-esimo raccoglie gli eventi caduti nel suo intervallo,
\begin{equation}
\label{eq:finestre_fisse}
\mathcal E_k=\bigl\{e_i\in\mathcal E\ \big|\ t_0+k\Delta t\leq t_i<t_0+(k+1)\Delta t\bigr\}.
\end{equation}
Con $t_0=0$ e $\Delta t=50\,\mathrm{ms}$, un evento all'istante $t_i=123\,\mathrm{ms}$ finisce nel bin $k=2$, quello che va da $100$ a $150\,\mathrm{ms}$. Il bordo sinistro è incluso e quello destro escluso, così che un evento che cade esattamente a $150\,\mathrm{ms}$ appartenga al bin successivo e mai a due bin insieme. In questo schema ogni passo dura lo stesso tempo, mentre il numero di eventi che contiene cambia liberamente da un passo all'altro.

Fissata la durata dei bin, resta da decidere quanti passi elaborare. La soluzione più comune è una finestra di durata prefissata, $D_{\mathrm{win}}=T\Delta t$, che produce sempre $T$ passi. Con bin da $50\,\mathrm{ms}$ e una finestra di $2\,\mathrm{s}$ i passi sono 40. Una registrazione di $1{,}2\,\mathrm{s}$ ne riempie i primi 24 e lascia vuoti gli ultimi 16, mentre di una registrazione di $2{,}5\,\mathrm{s}$ vanno persi gli ultimi $0{,}5\,\mathrm{s}$.

In alternativa si può conservare l'intera registrazione, dal primo evento, che fa da origine $t_0$, all'ultimo, che cade all'istante $t_{\mathrm{last}}$, mantenendo bin della stessa durata. Il numero di passi diventa allora
\begin{equation}
\label{eq:lunghezza_clock_fisico}
T=\left\lfloor\frac{t_{\mathrm{last}}-t_0}{\Delta t}\right\rfloor+1.
\end{equation}
La parte intera conta i bin completamente trascorsi prima dell'ultimo evento, e l'unità aggiunta fa posto al bin che lo contiene. Con bin da $50\,\mathrm{ms}$ e un ultimo evento a $1199\,\mathrm{ms}$ si ottengono 24 passi. Se l'ultimo evento cadesse esattamente a $1200\,\mathrm{ms}$, sul bordo sinistro di un nuovo bin, i passi sarebbero 25, ed è proprio l'unità aggiunta a garantire che quell'evento trovi posto. In questo schema la durata dei bin è la stessa per tutte le registrazioni, e a cambiare è la lunghezza della sequenza.

Una terza possibilità fa il contrario, fissa il numero di passi $T$ e adatta la durata dei bin a ciascuna registrazione $n$, di durata $D_n$,
\begin{equation}
\label{eq:delta_t_rapporto}
\Delta t_n=\frac{D_n}{T}.
\end{equation}
Con $T=40$, un gesto di $4\,\mathrm{s}$ viene diviso in bin da $100\,\mathrm{ms}$ e uno di $8\,\mathrm{s}$ in bin da $200\,\mathrm{ms}$. Il clock cambia così da un campione all'altro. Una distanza di cinque passi corrisponde a mezzo secondo nel primo gesto e a un secondo intero nel secondo, e la rete non ha modo di saperlo. Fissare il numero dei bin, in altre parole, significa non fissare la risoluzione temporale.

Si può infine rinunciare del tutto a misurare il tempo e contare gli eventi. Una prima variante chiude un gruppo ogni $M$ eventi, nell'ordine in cui arrivano. Una seconda divide l'intera registrazione in $T$ gruppi che contengono lo stesso numero di eventi, a meno di uno,
\begin{equation}
\label{eq:gruppi_eventi_bilanciati}
n_k\in\left\{\left\lfloor\frac NT\right\rfloor,
                 \left\lceil\frac NT\right\rceil\right\},
\qquad \sum_{k=1}^{T}n_k=N,
\end{equation}
dove $n_k$ è il numero di eventi del gruppo $k$. Con $N=10\,003$ eventi e $T=40$ gruppi, per esempio, tre gruppi ricevono 251 eventi e gli altri 37 ne ricevono 250. In questo schema ogni passo contiene la stessa quantità di attività, ma la sua durata dipende da quanto è intensa la scena. Una raffica concentrata in un decimo di secondo può riempire molti gruppi, mentre un secondo di quiete può starne in uno solo. La seconda variante richiede di conoscere in anticipo il numero totale $N$ degli eventi, e si può applicare soltanto a registrazioni già concluse, cioè offline. La prima non ha questo limite, perché ogni gruppo si chiude non appena raggiunge $M$ eventi. In entrambi i casi, però, una distanza misurata in gruppi non corrisponde a un intervallo di tempo fisso.

\begin{figure}[tbp]
    \centering
    \includegraphics[width=\textwidth]{images/cap1/cap1_partizioni.pdf}
    \caption[Quattro suddivisioni di uno stesso flusso di eventi]{Quattro modi di suddividere lo stesso flusso sintetico, che contiene una raffica di attività tra $0{,}35$ e $0{,}55\,\mathrm{s}$ e termina a circa $1{,}17\,\mathrm{s}$. In alto, la frequenza degli eventi. Sotto, i passi prodotti da ciascuna regola, colorati in base al numero di eventi che contengono. La finestra fissa lascia una coda di passi vuoti (tratteggiati), la registrazione intera termina con l'ultimo evento, il numero di passi fisso allarga i bin per coprire la durata, la partizione a eventi costanti produce passi con lo stesso contenuto ma di durata molto diversa, stretti durante la raffica e larghi altrove.}
    \label{fig:partizioni_flusso}
\end{figure}

Quando i passi seguono il numero di eventi, la loro durata non è più fissata in anticipo ma diventa una proprietà di ciascun gruppo, misurabile come intervallo (\emph{span}) tra il suo primo e il suo ultimo evento. Se la rete riceve soltanto i conteggi, però, questa durata le resta invisibile. Un gruppo di 250 eventi raccolti in $5\,\mathrm{ms}$ e uno raccolto in mezzo secondo occupano entrambi esattamente un passo.

La scelta del clock non è dunque un dettaglio di preprocessing. Da essa dipende il significato fisico di ogni ritardo espresso in passi e di ogni costante di decadimento della rete. Dimezzare $\Delta t$ su una stessa registrazione raddoppia il numero dei passi, dimezza in media i conteggi di ciascun bin e fa sì che un neurone che perde una frazione fissa del proprio potenziale a ogni passo dimentichi due volte più in fretta in tempo reale. I protocolli adottati in questo lavoro sono descritti nel capitolo~\ref{chap:baseline_metodologia} e messi a confronto nel capitolo~\ref{chap:clock_input}.

\subsection{Conteggi e informazione intra-bin}
\label{subsec:conteggi_intrabin}

Una volta suddiviso il flusso, bisogna decidere che cosa conservare di ciascun bin. La rappresentazione più semplice conta, per ogni pixel e per ogni bin, quanti eventi di ciascuna polarità vi sono caduti,
\begin{equation}
\label{eq:event_frame_accumulo}
F_k^p(x,y)=\sum_{e_i\in\mathcal E_k}
\mathbf 1[x_i=x,\ y_i=y,\ p_i=p],
\qquad F_k\in\mathbb N_0^{2\times H\times W}.
\end{equation}
L'indicatore $\mathbf 1[\cdot]$ vale 1 quando la condizione tra parentesi è vera e 0 altrimenti, quindi la somma conta semplicemente gli eventi del bin $k$ caduti nel pixel $(x,y)$ con polarità $p$. Il risultato è un'immagine a due canali, uno per gli eventi positivi e uno per quelli negativi, spesso chiamata \emph{event frame} (figura~\ref{fig:confronto_rappresentazioni_discrete}a). Se in un bin un pixel riceve tre eventi positivi e uno negativo, il canale positivo vale 3 e quello negativo 1.

Esistono versioni più povere della stessa idea. Una versione binaria, $\mathbf 1[F_k^p(x,y)>0]$, registra soltanto se il pixel è stato attivo, e nell'esempio darebbe 1 in entrambi i canali. Una versione a un solo canale somma gli eventi con il loro segno, e darebbe $3-1=2$. Quest'ultima ha un difetto evidente, perché due eventi positivi e due negativi nello stesso pixel si annullano e producono zero, come se non fosse accaduto nulla. Tenere separate le polarità evita il problema.

Tutte queste versioni condividono però un limite più sottile. Una somma non dipende dall'ordine degli addendi, e i conteggi di un bin non cambiano se gli eventi al suo interno vengono rimescolati. Si pensi a un oggetto chiaro che attraversa rapidamente un pixel scuro. Il pixel prima si illumina e poi torna scuro, ed emette un evento positivo seguito da uno negativo. Un oggetto scuro che attraversa un pixel chiaro produce la sequenza opposta. Se entrambi gli eventi cadono nello stesso bin, i due casi danno esattamente gli stessi conteggi, uno per canale, e nessuna rete, per quanto sofisticata, potrà distinguerli sulla base di quel solo bin. La rappresentazione fissa un limite all'informazione disponibile, che va tenuto distinto dalla capacità della rete di sfruttare quella che rimane.

La \emph{voxel grid} recupera parte dell'informazione temporale perduta~\cite{zhu2019unsupervised}. Si considera un pacchetto di eventi consecutivi, elaborato come un'unica unità, e se ne riscalano i timestamp in modo che il primo evento valga 0 e l'ultimo $B-1$,
\begin{equation}
\label{eq:timestamp_normalizzato}
t_i^*=(B-1)\frac{t_i-t_{\min}}{t_{\max}-t_{\min}},
\end{equation}
dove $B$ è il numero di bin temporali desiderati e il pacchetto deve avere durata non nulla, $t_{\max}>t_{\min}$. Ogni evento viene poi distribuito tra i due bin più vicini alla sua posizione $t_i^*$,
\begin{equation}
\label{eq:voxel_grid_interpolazione}
\mathcal V(b,x,y)=\sum_i p_i
\max\bigl(0,\,1-|b-t_i^*|\bigr)
\mathbf 1[x_i=x,\ y_i=y],
\qquad b\in\{0,\ldots,B-1\}.
\end{equation}
Il peso $\max(0, 1-|b-t_i^*|)$ vale 1 quando l'evento cade esattamente sul bin $b$, decresce in modo lineare con la distanza e si annulla a un bin di distanza. Con $B=5$, un evento positivo con $t_i^*=2{,}3$ aggiunge $0{,}7$ al bin 2 e $0{,}3$ al bin 3, e niente agli altri. L'evento non viene attribuito tutto a un bin o tutto all'altro, e la sua posizione temporale sopravvive, sfumata, nella ripartizione dei pesi (figura~\ref{fig:confronto_rappresentazioni_discrete}b). La formulazione completa interpola anche nello spazio, quando la correzione della distorsione ottica produce coordinate non intere~\cite{zhu2019unsupervised}.

Due dettagli vanno tenuti presenti. La somma usa la polarità con il suo segno, e può quindi produrre le stesse cancellazioni dell'event frame a un canale, che volumi separati per eventi positivi e negativi eviterebbero. Inoltre la scala temporale è relativa a ciascun pacchetto. Con $B=5$, un intervallo di $10\,\mathrm{ms}$ occupa meno di mezzo bin in un pacchetto lungo $100\,\mathrm{ms}$ e ben due bin in uno lungo $20\,\mathrm{ms}$, e non va confuso con un clock fisico comune a tutti i campioni.

La \emph{time surface} adotta un punto di vista diverso. Invece di contare gli eventi, descrive quanto sia recente l'attività di ciascun pixel. Nella forma più semplice,
\begin{equation}
\label{eq:time_surface}
S_p(x,y,t)=
\exp\left[-\frac{t-t_{\mathrm{last}}^p(x,y)}{\tau_s}\right],
\end{equation}
dove $t_{\mathrm{last}}^p(x,y)$ è l'istante dell'ultimo evento di polarità $p$ arrivato al pixel prima di $t$, e il valore è zero dove non è ancora arrivato nulla. Un pixel appena attivato vale circa 1, e il suo valore decade poi in modo esponenziale con costante $\tau_s$. Con $\tau_s=50\,\mathrm{ms}$, un pixel attivato $10\,\mathrm{ms}$ prima vale circa $0{,}82$, uno attivato $100\,\mathrm{ms}$ prima circa $0{,}14$. Il risultato somiglia alla scia di una cometa, intensa dove il movimento è appena passato e sempre più tenue alle sue spalle (figura~\ref{fig:confronto_rappresentazioni_discrete}c). È la variante elementare di una famiglia più ampia. HATS, per esempio, somma i contributi esponenziali di più eventi in un intorno spazio-temporale e li aggrega in descrittori~\cite{sironi2018hats}.

\begin{figure}[tbp]
    \centering
    \includegraphics[width=\textwidth]{images/cap1/cap1_rappresentazioni.pdf}
    \caption[Event frame, voxel grid e time surface a confronto]{Tre rappresentazioni degli stessi $100\,\mathrm{ms}$ di una registrazione di DVS128 Gesture, durante una rotazione del braccio sinistro (circa $9\,600$ eventi). (a)~Event frame, con i conteggi degli eventi positivi e negativi in due canali separati. (b)~Voxel grid con $B=5$ bin temporali, in cui la distribuzione dell'attività cambia da un bin all'altro, con i valori positivi in rosso e quelli negativi in blu. (c)~Time surface dell'ultimo evento di entrambe le polarità, valutata alla fine dell'intervallo, in cui la posizione più recente del movimento appare come una scia.}
    \label{fig:confronto_rappresentazioni_discrete}
\end{figure}

Event frame, voxel grid e time surface riassumono dunque aspetti diversi della stessa storia, la quantità dell'attività, la sua collocazione nel tempo e la sua recenza. La tabella~\ref{tab:confronto_rappresentazioni} ne riassume i compromessi. Una rappresentazione più ricca costa di più da costruire e da conservare, e la sua utilità dipende in ultima analisi dalle operazioni che il modello è in grado di applicarvi.

\begin{table}[tbp]
\centering
\small
\renewcommand{\arraystretch}{1.2}
\caption[Rappresentazioni degli eventi a confronto]{Informazione conservata e principali limiti delle rappresentazioni considerate.}
\label{tab:confronto_rappresentazioni}
\begin{tabularx}{\textwidth}{@{}>{\raggedright\arraybackslash}p{2.25cm} >{\raggedright\arraybackslash}X >{\raggedright\arraybackslash}X >{\raggedright\arraybackslash}X@{}}
\hline
\textbf{Formato} & \textbf{Informazione conservata} & \textbf{Descrizione del tempo} & \textbf{Limite principale} \\
\hline
Stream raw & Posizione, polarità e timestamp dei singoli eventi & Istanti originari & Sequenza irregolare e carico variabile \\
Conteggi ON/OFF & Molteplicità per pixel e polarità & Appartenenza al bin & Ordine intra-bin eliminato \\
Voxel grid & Accumuli su più bin temporali & Posizione temporale discretizzata o interpolata & Aggregazione e possibile cancellazione nella versione firmata \\
Time surface dell'ultimo evento & Recenza per pixel e polarità & Tempo trascorso dall'ultima attivazione & Molteplicità e storia anteriore non rappresentate separatamente \\
\hline
\end{tabularx}
\end{table}

\subsection{Tempo valido, silenzi e padding}
\label{subsec:tempo_valido_padding}

Quando le sequenze hanno lunghezze diverse sorge un problema pratico. Le reti vengono addestrate su batch, cioè gruppi di campioni elaborati insieme in un unico tensore, e un tensore richiede che tutte le sequenze abbiano la stessa lunghezza. Le sequenze più brevi vengono quindi allungate con passi fittizi, il cosiddetto \emph{padding}, fino alla lunghezza della più lunga, $T_{\max}$. Per tenere traccia di quali passi siano veri, a ogni campione $n$ di lunghezza $T_n$ si associa una maschera
\begin{equation}
\label{eq:maschera_passi_validi}
m_{k,n}=\mathbf 1[k<T_n],
\qquad k=0,\ldots,T_{\max}-1,
\end{equation}
che vale 1 sui passi osservati e 0 su quelli aggiunti. In un batch di tre sequenze lunghe 25, 33 e 40 passi, per esempio, tutte vengono portate a 40, e la maschera della prima vale 1 sui primi 25 passi e 0 sugli ultimi 15.

Un passo di padding e un bin valido senza eventi possono contenere entrambi soltanto zeri, ma significano cose opposte. Il bin vuoto dice che è trascorso del tempo senza attività, ed è a tutti gli effetti un'osservazione. Il padding non osserva nulla. Confonderli altera il campione. Se la rete calcolasse la media delle proprie feature su tutti i 40 passi anziché sui 25 validi, il risultato della prima sequenza verrebbe ridotto di oltre un terzo senza alcuna ragione legata ai dati. Lo stesso accade se il padding entra nelle statistiche della normalizzazione, e un modulo dotato di stato interno deve far avanzare il proprio tempo soltanto sui passi validi.

C'è infine una terza situazione, diversa da entrambe. Quando la registrazione finisce prima della finestra prevista, come nell'esempio della finestra di $2\,\mathrm{s}$, i passi successivi sono validi ma privi di eventi. La rete, però, non resta necessariamente inattiva, perché i potenziali di membrana e le memorie interne continuano a evolvere anche senza ingresso. Questa coda vuota e il padding di un batch vanno quindi dichiarati separatamente.

L'intera catena di elaborazione si può riassumere nello schema
\begin{equation}
\label{eq:decomposizione_flusso}
\mathcal E
\overset{\mathcal R}{\longrightarrow}X_{1:T}
\overset{\mathcal M}{\longrightarrow}Z_{1:T}
\overset{\mathcal D}{\longrightarrow}\hat y,
\end{equation}
dove $\mathcal R$ costruisce la rappresentazione, $\mathcal M$ ne elabora la dinamica e $\mathcal D$ prende la decisione. Ognuno dei tre stadi può limitare l'uso dell'informazione temporale. Per questo, quando si confrontano architetture diverse, rappresentazione e clock vanno tenuti fissi, e quando è l'ingresso stesso l'oggetto dell'esperimento ogni sua modifica va dichiarata come tale.


% ==============================================================================
%  Dinamica neuronale e memoria nelle Spiking Neural Networks
% ==============================================================================

\section{Dinamica neuronale e memoria nelle Spiking Neural Networks}
\label{sec:dinamica_neuronale_memoria}

Un neurone artificiale convenzionale, come quello di un percettrone multistrato con attivazione ReLU, calcola la propria uscita a partire dal solo ingresso corrente. Se riceve due volte lo stesso ingresso risponde due volte allo stesso modo, e non conserva nulla da un'elaborazione all'altra. Il neurone di una Spiking Neural Network (SNN) si comporta diversamente. Possiede una variabile interna, il potenziale di membrana, che accumula gli ingressi ricevuti e sopravvive da un passo al successivo, e la sua risposta dipende quindi anche da ciò che è accaduto prima. È la struttura di una rete ricorrente, con una differenza importante. In una rete ricorrente classica il passato rientra attraverso connessioni esplicite che riportano l'uscita di uno strato al suo ingresso, mentre in una SNN la ricorrenza è già dentro ogni neurone, e la rete ha memoria anche quando i suoi strati sono collegati soltanto in avanti~\cite{neftci2019surrogate}.

Questa sezione descrive come un neurone accumuli e dimentichi gli ingressi, come la velocità con cui dimentica dipenda dal clock dell'input, come si addestri una rete che comunica con impulsi discontinui e che cosa una memoria di questo tipo permetta di fare.

\subsection{Integrazione, soglia e reset}
\label{subsec:integrazione_soglia_reset}

Il modello più semplice è l'Integrate-and-Fire (IF), che si può immaginare come un accumulatore. A ogni passo il neurone somma al proprio potenziale il contributo dell'ingresso,
\begin{equation}
\label{eq:if_discreto}
H_k=V_{k-1}+X_k,
\end{equation}
e quando il potenziale raggiunge una soglia $V_{\mathrm{th}}$ emette uno spike,
\begin{equation}
\label{eq:if_spike}
S_k=\Theta(H_k-V_{\mathrm{th}}),
\qquad
\Theta(z)=
\begin{cases}
1,&z\geq0,\\
0,&z<0.
\end{cases}
\end{equation}
La funzione a gradino $\Theta$ traduce il confronto con la soglia in una risposta tutto o niente. Lo spike vale 1 se il potenziale ha raggiunto la soglia e 0 altrimenti, senza valori intermedi, e un potenziale appena sopra la soglia produce lo stesso spike di uno che la supera di molto.

Le due equazioni usano due simboli diversi per il potenziale, e conviene chiarire subito perché. All'interno di ogni passo avvengono due cose in sequenza. Prima il neurone riceve l'ingresso e si carica, e $H_k$ indica il potenziale in questo momento, prima che il neurone decida se emettere lo spike. Poi, se lo spike è stato emesso, il potenziale viene riportato verso il basso, ed è il valore che resta dopo questa operazione, indicato con $V_k$, a passare al passo successivo. Se lo spike non viene emesso, i due valori coincidono.

L'operazione che segue lo spike si chiama reset, e ne esistono due varianti. L'\emph{hard reset} riporta il potenziale a un valore fisso, di solito zero, mentre il \emph{soft reset}, adottato in questo lavoro, si limita a sottrarre la soglia,
\begin{equation}
\label{eq:soft_reset}
V_k=H_k-S_k V_{\mathrm{th}}.
\end{equation}
La differenza emerge quando il potenziale supera la soglia di un certo margine. Con soglia 1 e $H_k=1{,}3$, l'hard reset riparte da zero e perde l'eccedenza di $0{,}3$, mentre il soft reset riparte da $0{,}3$ e la conserva per il passo successivo~\cite{han2020rmp}. Qui e nel seguito si adotta la convenzione con potenziale di riposo nullo.

Il neurone IF ha una memoria perfetta. Al di fuori del reset, un contributo ricevuto all'inizio della sequenza pesa esattamente quanto uno appena arrivato, perché il modello non prevede alcuna forma di oblio.\\
Il Leaky Integrate-and-Fire (LIF) introduce proprio questo elemento, una perdita o \emph{leak}. A ogni passo, prima di sommare il nuovo ingresso, il potenziale perde una frazione fissa del proprio valore,
\begin{equation}
\label{eq:lif_decadimento_discreto}
H_k=\lambda V_{k-1}+\eta X_k,
\qquad 0<\lambda<1.
\end{equation}
Il fattore di decadimento $\lambda$ indica quale frazione del potenziale sopravvive da un passo al successivo, e il coefficiente $\eta$ stabilisce quanto pesi l'ingresso. Con $\lambda=0{,}5$, in assenza di ingresso, il potenziale si dimezza a ogni passo. Il neurone IF si ritrova ponendo $\lambda=1$ ed $\eta=1$.

La perdita cambia in modo sostanziale il comportamento del neurone, come mostra la figura~\ref{fig:dinamica_lif}. Un neurone con soglia 1 riceve tre ingressi da $0{,}6$, una volta in tre passi consecutivi e una volta a tre passi di distanza l'uno dall'altro. Il neurone IF emette uno spike in entrambi i casi, non appena il secondo ingresso porta il potenziale a $1{,}2$. Il neurone LIF con $\lambda=0{,}5$ ed $\eta=1$, quando gli ingressi sono ravvicinati, sale a $0{,}6$, poi a $0{,}9$ e infine a $1{,}05$, e supera la soglia al terzo passo. Quando invece sono distanziati, tra un ingresso e l'altro il potenziale ha il tempo di scendere, non va mai oltre $0{,}68$ e lo spike non arriva. A parità di ingresso complessivo, il LIF risponde agli ingressi che arrivano vicini e lascia cadere quelli isolati, e si comporta in questo senso come un rilevatore di coincidenze.

\begin{figure}[t]
    \centering
    \includegraphics[width=\textwidth]{images/cap1/cap1_dinamica_lif.pdf}
    \caption[Integrazione, soglia e soft reset nei neuroni IF e LIF]{Integrazione, soglia e soft reset in un neurone IF (tratteggio arancione) e in un neurone LIF con $\lambda=0{,}5$ (linea blu), entrambi con soglia $V_{\mathrm{th}}=1$. I due neuroni ricevono tre ingressi da $0{,}6$ (in basso), consecutivi in (a) e distanziati in (b). I marcatori indicano il potenziale $H_k$ appena caricato, i tratti verticali il soft reset che segue uno spike, la fila in alto gli spike emessi. Il neurone IF, che non dimentica, emette uno spike in entrambi i casi. Il LIF lo emette soltanto quando gli ingressi sono abbastanza ravvicinati da sommarsi prima che il potenziale decada.}
    \label{fig:dinamica_lif}
\end{figure}

\subsection{Decadimento discreto e scala temporale fisica}
\label{subsec:decadimento_scala_fisica}

Il fattore $\lambda$ è definito per passo, mentre la memoria di un neurone, dal punto di vista della scena osservata, si misura in millisecondi. Per collegare le due cose bisogna risalire al modello da cui il LIF discende. Nella formulazione originale, a tempo continuo, la membrana si comporta come un condensatore che si scarica lentamente attraverso una resistenza,
\begin{equation}
\label{eq:lif_continuo}
\tau_m\frac{\mathrm d V(t)}{\mathrm dt}
=-\bigl(V(t)-V_{\mathrm{rest}}\bigr)+R I(t),
\end{equation}
dove $I(t)$ è la corrente d'ingresso, $R$ la resistenza di membrana, $V_{\mathrm{rest}}$ il potenziale di riposo e $\tau_m$ la costante di tempo di membrana~\cite{gerstner2014neuronal}. Il primo termine a destra riporta il potenziale verso il riposo con una spinta proporzionale alla distanza, il secondo lo sposta in base all'ingresso. Se l'ingresso si annulla, il potenziale torna verso il riposo lungo un'esponenziale. Trascorso un tempo pari a $\tau_m$ la distanza dal riposo si è ridotta a circa il $37\%$, poiché $e^{-1}\approx0{,}37$, e la costante $\tau_m$ misura dunque, in unità di tempo, per quanto il neurone ricorda un ingresso.

Per simulare questa dinamica a passi di durata $\Delta t$ occorre discretizzarla. Molte implementazioni, tra cui la libreria SpikingJelly~\cite{fang2023spikingjelly}, adottano la discretizzazione di Eulero, che sostituisce la derivata con un rapporto incrementale e ritrova proprio la forma dell'equazione~\eqref{eq:lif_decadimento_discreto},
\begin{equation}
\label{eq:lif_euler_discreto}
H_k=V_{k-1}+\frac{X_k-V_{k-1}}{\tau_d},
\qquad
\lambda=1-\frac1{\tau_d},\quad
\eta=\frac1{\tau_d}.
\end{equation}
Qui la costante $\tau_d$, che deve superare 1 perché $\lambda$ resti positivo, non è un tempo ma un numero puro, espresso in passi. Letta come discretizzazione dell'equazione continua vale $\tau_d=\tau_m/\Delta t$, ma Eulero approssima soltanto la soluzione esatta, per la quale varrebbe $\lambda=e^{-\Delta t/\tau_m}$. Con una costante pari a due passi si ottiene $\lambda=0{,}5$ nel primo caso e circa $0{,}61$ nel secondo. Quando un lavoro riporta una costante di tempo, quindi, bisogna sapere se si tratti di un tempo fisico o di un numero di passi, e con quale discretizzazione venga usata.

Srotolando la ricorrenza, cioè sostituendo a ritroso ogni potenziale con la sua espressione, si vede che cosa il neurone conservi della propria storia. In assenza di spike e reset, con coefficienti costanti, si ottiene
\begin{equation}
\label{eq:lif_espansione_temporale}
V_k=\lambda^k V_0+
\eta\sum_{j=1}^{k}\lambda^{k-j}X_j.
\end{equation}
Il potenziale è la somma di tutti gli ingressi ricevuti, ciascuno moltiplicato per $\lambda$ tante volte quanti sono i passi trascorsi dal suo arrivo. Con $\lambda=0{,}5$, rispetto all'ingresso appena arrivato, quello del passo precedente pesa la metà, quello di tre passi prima un ottavo e quello di dieci passi prima meno di un millesimo. Gli ingressi passati, però, non vi compaiono come osservazioni separate. Sono fusi in un unico numero, e su questa proprietà si tornerà alla fine della sezione.

La stessa espressione dice come tradurre $\lambda$ in un tempo fisico. Un ingresso arrivato $n$ passi prima, cioè $n\Delta t$ secondi prima, pesa $\lambda^n$. Cercando la costante $\tau_{\mathrm{eq}}$ per cui $\lambda^n=\exp(-n\Delta t/\tau_{\mathrm{eq}})$, come nel rilassamento continuo, si trova
\begin{equation}
\label{eq:tau_fisica_equivalente}
\tau_{\mathrm{eq}}=-\frac{\Delta t}{\ln\lambda}.
\end{equation}
Con $\lambda=0{,}5$ vale circa $1{,}44\,\Delta t$, quindi poco più di $14\,\mathrm{ms}$ con bin da $10\,\mathrm{ms}$ e poco più di $140\,\mathrm{ms}$ con bin da $100\,\mathrm{ms}$.

Il punto essenziale è che $\tau_{\mathrm{eq}}$ dipende dalla durata del bin, cioè dal clock dell'input. Lo stesso neurone, con lo stesso $\lambda$, dimentica due volte più in fretta in tempo reale se la durata dei bin si dimezza. Adeguare $\lambda$ al clock conserverebbe la memoria fisica, ma non renderebbe equivalenti i due modelli, perché cambiando i bin cambiano anche i conteggi che ciascuno contiene e il numero di occasioni in cui il neurone può emettere spike. Con bin di durata variabile, come nelle partizioni a eventi costanti, un coefficiente fisso per passo non corrisponde infine ad alcun decadimento uniforme nel tempo fisico, perché un passo lungo mezzo secondo e uno lungo cinque millisecondi fanno perdere al potenziale la stessa frazione.

Queste relazioni descrivono un singolo neurone isolato, al di sotto della soglia. La memoria di una rete, che combina molti neuroni, trasformazioni apprese e interazioni tra strati, non ha un orizzonte massimo così rigidamente definito.

\subsection{Addestramento con surrogate gradient}
\label{subsec:surrogate_gradient}

Le reti neurali si addestrano con la retropropagazione dell'errore, che calcola per ogni peso di quanto cambierebbe la loss modificandolo di poco, e lo corregge di conseguenza. La soglia che genera gli spike rende questo calcolo inutilizzabile. La funzione $\Theta$ è piatta ovunque tranne che nell'origine, dove salta da 0 a 1. Se il potenziale cambia di poco lo spike resta quello che era, e la derivata è quindi nulla quasi ovunque, mentre nell'unico punto in cui lo spike cambia la derivata non è definita. In entrambi i casi nessun segnale di errore riesce ad attraversare la soglia e a raggiungere i pesi degli strati precedenti.

I \emph{surrogate gradient} aggirano l'ostacolo con un compromesso~\cite{neftci2019surrogate}. Nel passaggio in avanti (forward pass), quello che produce le uscite della rete, lo spike resta esattamente quello dell'equazione~\eqref{eq:if_spike}. Nel passaggio all'indietro (backward pass), quello che calcola i gradienti, la sua derivata viene sostituita da una funzione regolare $g_\alpha(H_k-V_{\mathrm{th}})$. Una scelta diffusa è la derivata di una sigmoide. Con $\sigma(u)=(1+e^{-u})^{-1}$ si pone
\begin{equation}
\label{eq:surrogate_sigmoide}
g_\alpha(z)=\alpha\,\sigma(\alpha z)
                    \bigl[1-\sigma(\alpha z)\bigr],
\qquad \alpha>0,
\end{equation}
che è proprio la derivata di $\sigma(\alpha z)$. In pratica la rete calcola gli spike con il gradino ma, al momento di calcolare i gradienti, si comporta come se il gradino fosse stato una sigmoide (figura~\ref{fig:surrogate_gradient_schema}). Il parametro $\alpha$ regola la ripidità della sigmoide, e quindi l'altezza e la larghezza del gradiente surrogato, che raggiunge il massimo sulla soglia e si annulla rapidamente allontanandosene. Ricevono così un gradiente apprezzabile soltanto i neuroni il cui potenziale era vicino alla soglia, quelli cioè per cui una piccola modifica dei pesi avrebbe potuto cambiare la decisione di emettere o meno lo spike.

\begin{figure}[t]
    \centering
    \includegraphics[width=\textwidth]{images/cap1/cap1_surrogata.pdf}
    \caption[Principio del surrogate gradient]{Principio del surrogate gradient, in funzione della distanza $z=H_k-V_{\mathrm{th}}$ tra potenziale e soglia. (a)~Nel passaggio in avanti lo spike è prodotto dal gradino $\Theta$ (linea nera). La sigmoide $\sigma(4z)$ (tratteggio) non viene mai calcolata, ma è la funzione di cui il passaggio all'indietro usa la derivata. (b)~La derivata del gradino è nulla ovunque tranne nell'origine, dove non è definita (freccia), e non trasmette alcun gradiente. La surrogata $g_\alpha$ la sostituisce con una campana centrata sulla soglia, tanto più alta e stretta quanto maggiore è $\alpha$. Il valore evidenziato, $\alpha=4$, è quello adottato in questo lavoro.}
    \label{fig:surrogate_gradient_schema}
\end{figure}

La soglia non è l'unico ostacolo. Poiché il potenziale collega ogni passo al successivo, la loss dipende dai pesi anche attraverso il tempo. La Backpropagation Through Time (BPTT) affronta il problema srotolando la ricorrenza~\cite{wu2018spatiotemporal}. La rete viene vista come una catena di $T$ copie di sé stessa, una per ciascun passo, che condividono gli stessi pesi e sono collegate dagli stati che ciascuna passa alla successiva. Su questa catena, che non contiene più cicli, la retropropagazione ordinaria si applica senza difficoltà, e l'errore si propaga in due direzioni, attraverso gli strati e all'indietro nel tempo. Passando da un passo al precedente attraverso la membrana, il gradiente viene moltiplicato all'incirca per $\lambda$, lo stesso fattore con cui il neurone dimentica. Un neurone che dimentica in fretta impara quindi poco dalle conseguenze lontane dei propri stati, una forma della scomparsa del gradiente già nota nelle reti ricorrenti~\cite{pascanu2013difficulty}.

Il prezzo della BPTT è la memoria. Per percorrere la catena all'indietro servono i valori calcolati in avanti in ogni passo, e il solo stato neuronale da conservare cresce come $\mathcal O(T\sum_l N_l)$, dove $N_l$ è il numero di unità dello strato $l$~\cite{neftci2019surrogate}. In inferenza, invece, a una rete di soli neuroni LIF basta il potenziale attuale di ciascun neurone, cioè $\sum_l N_l$ valori qualunque sia $T$. Le due grandezze vanno tenute distinte, perché un modello può richiedere pochissimo stato durante l'esecuzione e, al tempo stesso, un grafo molto esteso durante l'ottimizzazione.

In questo lavoro la rete spiking viene addestrata direttamente, con surrogate gradient e BPTT. È un paradigma diverso dalla conversione ANN-SNN, in cui una rete convenzionale viene prima addestrata e poi tradotta in una rete spiking che ne approssima le attivazioni con la frequenza degli spike. La dinamica discreta, la funzione surrogata e il trattamento del reset effettivamente adottati sono specificati nella sezione~\ref{subsec:neurone_convenzioni}.

\subsection{Memoria neuronale e relazioni tra feature}
\label{subsec:memoria_relazioni}

Si torni all'equazione~\eqref{eq:lif_espansione_temporale}. Per un neurone LIF privo di ulteriori stati interni, tutto ciò che la storia degli ingressi può trasmettere al futuro passa per un unico numero, il potenziale $V_k$. L'intera sequenza viene compressa in un solo valore, e storie diverse possono lasciare lo stesso potenziale. Con $\lambda=0{,}5$ ed $\eta=1$, partendo da un potenziale nullo, un ingresso pari a $0{,}4$ arrivato al passo precedente e un ingresso pari a $0{,}2$ arrivato adesso lasciano entrambi $V_k=0{,}2$. Da quel momento, se gli ingressi successivi coincidono, il neurone si comporta in modo identico nei due casi, e non ha più modo di sapere se abbia visto qualcosa di intenso poco fa o qualcosa di debole in questo istante. La rete nel suo complesso può comunque conservare la differenza in altre unità o in altri moduli, sicché la compressione operata da un singolo potenziale non implica di per sé una perdita equivalente a livello dell'intera rete.

Si delineano così due capacità diverse. La prima è la \emph{persistenza}, cioè la capacità di conservare a lungo una traccia del passato. La seconda è l'\emph{accesso}, cioè la possibilità di recuperare separatamente ciò che è accaduto in un istante preciso per confrontarlo con il presente. Si potrebbe pensare di migliorare la memoria di un neurone avvicinando $\lambda$ a 1, ma la modifica agisce soltanto sulla prima. I contributi passati sopravvivono più a lungo, e proprio per questo un numero maggiore di ingressi si mescola nella stessa somma, che diventa ancora più difficile da scomporre. Si pensi alla barra della sezione~\ref{subsec:contrasto_eventi}, la cui direzione è scritta nell'ordine con cui si attivano pixel vicini. Per riconoscerla non basta sapere che entrambi i pixel sono stati attivi di recente, informazione che un potenziale conserva bene. Bisogna sapere quale dei due si è attivato prima, cioè mettere a confronto due istanti distinti, e un'unica somma pesata può farlo soltanto in modo indiretto.

La distinzione tra persistenza e accesso motiva l'interesse per operatori capaci di combinare o confrontare esplicitamente feature diverse. La sezione successiva descrive come l'attenzione sia stata integrata nelle architetture spiking, e sarà allora necessario distinguere le relazioni costruite all'interno di uno stesso passo da quelle che collegano istanti diversi.


% ==============================================================================
%  Dalle SNN agli Spiking Transformer
% ==============================================================================


\section{Dalle SNN agli Spiking Transformer}
\label{sec:snn_spiking_transformer}

La dinamica neuronale offre alle SNN un modo naturale di trasportare informazione da un passo temporale, o \emph{timestep}, al successivo. Resta però aperto un problema di altra natura, quello di mettere in relazione le numerose feature che compongono una rappresentazione visiva. Un neurone LIF integra nel tempo il proprio ingresso, ma nulla nella sua dinamica gli permette di stabilire quali regioni della scena siano rilevanti l'una per l'altra, né di collegare feature spazialmente lontane. Questa sezione ripercorre il modo in cui l'attenzione dei Transformer, nata per affrontare questo problema nelle reti convenzionali, è stata adattata al dominio spiking, fino all'architettura da cui muove la tesi.

\subsection{Self-attention e relazioni tra token}
\label{subsec:self_attention}

Le architetture convoluzionali affrontano il problema componendo operatori locali. Ogni convoluzione combina l'informazione di un intorno limitato, e soltanto impilando molti strati, che allargano progressivamente il campo recettivo, la rete arriva a costruire relazioni spaziali estese. Su questo principio si fonda l'efficacia delle CNN nella visione artificiale, ed è su di esso che sono state costruite molte delle prime SNN profonde.

Il Transformer adotta una prospettiva diversa. Anziché affidare la costruzione delle dipendenze a una lunga catena di operatori locali, permette agli elementi di una rappresentazione di interagire direttamente attraverso il meccanismo di self-attention~\cite{vaswani2017attention}. Trasferita alla visione, l'idea consiste nel suddividere l'immagine in porzioni, le patch, ciascuna delle quali viene proiettata in un vettore di feature detto token~\cite{dosovitskiy2021image}. I token possono poi influenzarsi a vicenda in base al contenuto che rappresentano, qualunque sia la distanza che li separa nell'immagine.

Il meccanismo si capisce bene a partire da un dizionario, la struttura dati che associa a ogni chiave un valore. Per recuperare un valore si presenta una richiesta (\emph{query}), la si confronta con le chiavi (\emph{key}) e si ottiene il valore (\emph{value}) associato alla chiave identica alla richiesta. Un dizionario di traduzioni, interrogato con la parola «gatto», trova la chiave «gatto» e restituisce «cat». L'attenzione rende graduale questa ricerca (figura~\ref{fig:attenzione_dizionario}). Richiesta e chiavi diventano vettori di feature, e il confronto non dà più un esito sì o no ma un punteggio di somiglianza con ciascuna chiave. Il risultato, di conseguenza, non è un singolo valore ma una media di tutti i valori, pesata con quei punteggi. Con pesi pari a $0{,}7$, $0{,}2$ e $0{,}1$ si ottiene $0{,}7\,v_1+0{,}2\,v_2+0{,}1\,v_3$, dove conta soprattutto il valore associato alla chiave più affine alla richiesta.

\begin{figure}[t]
    \centering
    \includegraphics[width=\textwidth]{images/cap1/cap1_attenzione_dizionario.pdf}
    \caption[L'attenzione come ricerca graduale in un dizionario]{L'attenzione come ricerca graduale in un dizionario. (a)~Un dizionario confronta la richiesta con le chiavi e restituisce soltanto il valore associato alla chiave identica. (b)~L'attenzione confronta la Query con tutte le Key, ricava da ciascun confronto un peso e restituisce la media dei Value pesata in questo modo. Lo spessore dei collegamenti è proporzionale al peso, e i valori numerici sono illustrativi.}
    \label{fig:attenzione_dizionario}
\end{figure}

Nella self-attention richiesta, chiavi e valori provengono tutti dagli stessi token, ed è da qui che viene il prefisso \emph{self}. Data una matrice di token $X\in\mathbb{R}^{N\times D}$, con $N$ token di dimensione $D$, tre trasformazioni lineari apprese ricavano da ciascun token tre vettori,
\begin{equation}
\label{eq:qkv_projections}
Q = X W_Q, \qquad K = X W_K, \qquad V = X W_V,
\end{equation}
la Query, la Key e il Value. Ogni token ricopre quindi tre ruoli. Con la Query esprime che cosa cerca negli altri token, con la Key descrive che cosa offre a chi lo interroga, con il Value fornisce il contenuto che gli altri riceveranno se lo trovano rilevante. Tenere separati i tre vettori permette alla rete di imparare che ciò che un token cerca può essere diverso da ciò che offre. Il token che copre un tratto di contorno, per esempio, può imparare a cercare i token che ne proseguono il profilo, anche se si trovano dall'altra parte dell'immagine.

Le formule traducono direttamente questa ricerca. L'elemento $(i,j)$ del prodotto $QK^{\top}$ è il prodotto scalare tra la Query del token $i$ e la Key del token $j$, cioè il punteggio di somiglianza tra ciò che $i$ cerca e ciò che $j$ offre. La softmax trasforma ogni riga di punteggi in pesi positivi che sommano a 1, e questi pesi combinano i Value,
\begin{equation}
\label{eq:vanilla_softmax_attention}
Y = \operatorname{softmax}\left( \frac{Q K^{\top}}{\sqrt{d}} \right) V,
\end{equation}
dove la divisione per $\sqrt{d}$, con $d$ dimensione di Query e Key, evita che i prodotti scalari crescano troppo con la dimensione. La riga $i$ di $Y$ è il nuovo token $i$, cioè la media dei Value di tutti i token pesata secondo la loro compatibilità con la sua Query. Nel seguito questa forma viene indicata come VSA (\emph{Vanilla Self-Attention}). Ne deriva una capacità che la convoluzione non possiede nella stessa forma. Due regioni lontane possono interagire direttamente all'interno di un singolo blocco, senza che l'informazione debba attraversare una lunga successione di operatori locali.

Questa libertà ha un prezzo. Poiché $QK^{\top}$ confronta ogni token con tutti gli altri, la matrice di attenzione ha dimensione $N\times N$, il costo del prodotto cresce come $\mathcal{O}(N^2D)$ e la memoria necessaria a conservare la matrice come $\mathcal{O}(N^2)$. Con 256 token la matrice contiene già più di $65\,000$ elementi. In una SNN se ne aggiunge un secondo vincolo, perché le stesse operazioni vanno ripetute a ogni timestep. Per un ingresso $X\in\mathbb{R}^{T\times N\times D}$ l'attenzione viene tipicamente applicata, istante per istante, alle feature di ciascun passo, e la dimensione temporale, che finora era apparsa come la sorgente della memoria neuronale, diventa così anche un moltiplicatore del costo.

\subsection{La self-attention in forma spiking}
\label{subsec:ssa}

L'incontro tra SNN e Transformer nasce da questa tensione. La VSA si basa su attivazioni continue, prodotti tra matrici dense e normalizzazioni, operazioni che non traggono alcun vantaggio dalla natura binaria e sparsa degli spike. Sostituire le attivazioni di un Transformer con neuroni LIF non basta dunque a ottenere uno Spiking Transformer. Occorre ripensare almeno in parte il meccanismo di attenzione, altrimenti la rappresentazione spiking verrebbe subito riconvertita in una computazione del tutto analoga a quella di una rete convenzionale (ANN).

Per capire che cosa si voglia preservare conviene anticipare un'idea che la sezione~\ref{sec:efficienza_neuromorfica_vincoli} riprenderà. Per applicare un peso a un'attivazione continua serve una moltiplicazione seguita da un accumulo (MAC). Se l'attivazione è uno spike binario la moltiplicazione scompare, perché basta sommare il peso quando lo spike c'è e non fare nulla quando manca, con un semplice accumulo (AC). Una rete si dice \emph{spike-driven} quando funziona così, cioè quando ogni operazione è innescata da uno spike e l'assenza di uno spike coincide con l'assenza di calcolo~\cite{yao2023spikedriven}.

Con Spikformer, Zhou et al. propongono nel 2023 la Spiking Self-Attention (SSA), mostrando che Query, Key e Value possono essere espressi in forma spiking e che la softmax può essere eliminata~\cite{zhou2023spikformer}. Dalle feature spiking si ricavano tre rappresentazioni binarie,
\begin{equation}
\label{eq:spiking_qkv}
Q_s, \; K_s, \; V_s \in \{0, 1\}^{N \times D},
\end{equation}
la cui interazione continua a svolgere il ruolo dell'attenzione. In forma semplificata,
\begin{equation}
\label{eq:ssa_spikformer}
\operatorname{SSA}(Q_s, K_s, V_s) = \operatorname{SN}\left( Q_s K_s^{\top} V_s \, \beta \right),
\end{equation}
dove $\operatorname{SN}(\cdot)$ indica lo strato di neuroni spiking che riporta il risultato in forma binaria e $\beta$ è un fattore di scala. La rinuncia alla softmax ha un peso particolare, per due ragioni. La softmax richiede esponenziali, divisioni e una normalizzazione che coinvolge tutti i coefficienti di una riga, operazioni estranee a una computazione guidata dagli eventi. Nel dominio spiking, inoltre, viene meno il suo compito principale, quello di rendere positivi i pesi, perché il prodotto di due matrici binarie non può produrre valori negativi~\cite{zhou2023spikformer}.

Spikformer dimostra così che self-attention e SNN sono compatibili, ma porta alla luce due questioni. La prima riguarda la natura effettiva delle operazioni. Che Query, Key e Value siano spike non garantisce che l'intero blocco esegua soltanto operazioni economiche. L'elemento $(i,j)$ di $Q_sK_s^{\top}$ conta in quanti canali la Query del token $i$ e la Key del token $j$ sono attive insieme. Se la prima vale $[1,0,1,1]$ e la seconda $[1,1,1,0]$ il risultato è 2, perché coincidono nel primo e nel terzo canale. Calcolare questo numero è economico, perché basta incrementare un contatore quando entrambi i bit valgono 1. Il risultato, però, non è più uno spike ma un intero, che può arrivare al numero dei canali, e lo strato SN lo riporta in forma binaria soltanto dopo che è stato combinato con i Value e scalato. Un problema analogo si ripresenta nelle connessioni residue, dove la somma di due spike può già valere 2. Quando un valore intero raggiunge uno strato dotato di pesi, ciascun peso va moltiplicato per quel valore, e ricompaiono le moltiplicazioni che il paradigma spike-driven vorrebbe evitare.

La seconda questione riguarda la scalabilità. Il prodotto $Q_sK_s^{\top}$ conserva una matrice tra tutte le coppie di token, e anche con dati sparsi la dimensione logica dell'operazione resta legata a $N^2$, un costo che diventa critico al crescere del numero dei token. Da questi due problemi sono nate linee di ricerca diverse e complementari, che le prossime sottosezioni esaminano nell'ordine.

\subsection{Attenzione e residui spike-driven}
\label{subsec:attenzione_residui}

Spike-driven Transformer affronta la prima questione, chiedendosi quanta parte della rete sia davvero spike-driven~\cite{yao2023spikedriven}. Gli autori osservano che diverse architetture precedenti mescolano accumuli e moltiplicazioni, dando luogo a una computazione ibrida, e si pongono l'obiettivo di mantenere binaria la comunicazione tra i blocchi. Il loro meccanismo, la Spike-Driven Self-Attention (SDSA), abbandona il prodotto completo $Q_sK_s^{\top}$. In una delle sue formulazioni, Query e Key vengono moltiplicate elemento per elemento e sommate lungo la dimensione dei token, e il risultato attraversa uno strato spiking,
\begin{equation}
\label{eq:sdsa_vector}
A_s = \operatorname{SN}\left( \operatorname{SUM}_N \left( Q_s \odot K_s \right) \right),
\end{equation}
dopodiché il vettore binario così ottenuto viene applicato ai Value,
\begin{equation}
\label{eq:sdsa_value}
Y_s = A_s \odot V_s,
\end{equation}
dove $\odot$ indica il prodotto elemento per elemento. Con operandi binari questo prodotto equivale a una maschera, perché ogni elemento viene semplicemente conservato oppure azzerato. Il procedimento si legge in tre tempi. Il prodotto $Q_s\odot K_s$ segna dove Query e Key sono attive insieme, la somma conta per ciascun canale quante volte ciò accade tra tutti i token, e il neurone trasforma questi conteggi in una decisione binaria, canale per canale. I canali promossi attraversano il blocco, gli altri vengono azzerati.

\begin{figure}[t]
    \centering
    \includegraphics[width=\textwidth]{images/cap1/cap1_operatori_attenzione.pdf}
    \caption[Riformulazione della self-attention nel dominio spiking]{Riformulazione della self-attention nel dominio spiking, calcolata su un piccolo esempio casuale con 6 token (righe) e 5 canali (colonne). Il colore indica il tipo di valore, e la cornice spessa la struttura che svolge il ruolo di attenzione. Dalla VSA alla SSA scompare la softmax e i dati in ingresso diventano binari, ma la matrice tra coppie di token resta e contiene conteggi. SDSA e QKTA sostituiscono la matrice con un vettore binario, che in uscita azzera interi canali nel primo caso e interi token nel secondo. La soglia degli strati SN è scelta per rendere leggibile l'esempio.}
    \label{fig:evoluzione_self_attention_spiking}
\end{figure}

Rispetto all'equazione~\eqref{eq:vanilla_softmax_attention} il cambiamento è sostanziale, come mostra la figura~\ref{fig:evoluzione_self_attention_spiking}. Non viene costruita alcuna matrice $N\times N$ di relazioni tra coppie di token, e l'attenzione si riduce a una scelta su quali canali lasciar passare. Per questo SDSA può essere interpretata come una forma particolare di attenzione lineare, la cui complessità cresce linearmente con il numero di token e di canali anziché quadraticamente con il numero di token~\cite{yao2023spikedriven}. Si rinuncia così a parte della ricchezza delle correlazioni tra coppie, in cambio di un'operazione molto più adatta alla computazione sparsa.

Spike-driven Transformer interviene anche su un elemento meno appariscente, ma importante per comprendere gli sviluppi successivi, cioè la posizione delle connessioni residue. Una connessione residua somma l'ingresso di un blocco alla trasformazione che il blocco stesso produce, e nelle reti profonde facilita la propagazione dell'informazione e dei gradienti. In una SNN, però, conta dove la somma viene eseguita, se tra spike oppure tra potenziali di membrana (figura~\ref{fig:residui_sew_ms}).

Nella connessione residua Spike-Element-Wise (SEW) adottata da Spikformer si sommano direttamente due tensori di spike. Il risultato può valere $0$, $1$ o $2$, e l'intervallo si allarga a ogni blocco attraversato, perché la somma di un blocco diventa un addendo di quella del successivo. Il segnale smette così di essere binario. Quando raggiunge uno strato pesato non basta più selezionare i pesi corrispondenti agli spike attivi, perché ciascun peso va moltiplicato per un valore intero, e ricompaiono moltiplicazioni tra interi e numeri in virgola mobile.

\begin{figure}[t]
    \centering
    \includegraphics[width=\textwidth]{images/cap1/cap1_residui.pdf}
    \caption[Connessioni residue SEW e MS]{Due posizioni della connessione residua, indicate con i valori che ogni tratto può assumere~\cite{zhou2026spikingformer}. (a)~Nel residuo SEW la somma avviene tra spike. I valori crescono a ogni blocco, e lo strato pesato successivo deve moltiplicare ciascun peso per un intero. (b)~Nel residuo MS la somma avviene tra potenziali di membrana, cioè tra valori reali, e un neurone spiking precede ogni strato pesato, che riceve così soltanto ingressi binari.}
    \label{fig:residui_sew_ms}
\end{figure}

Spike-driven Transformer sposta invece la somma residua sui potenziali di membrana, e colloca un neurone spiking subito prima di ogni strato pesato. La somma avviene così tra valori reali, come quella che ogni neurone esegue già al proprio interno, e ciò che entra negli strati pesati resta binario~\cite{yao2023spikedriven}. Il dettaglio contiene una lezione progettuale di portata più generale. Contare i neuroni spiking presenti in un'architettura dice poco sulla sua natura event-driven, che dipende piuttosto dalla posizione di quei neuroni rispetto alle operazioni lineari e ai percorsi residui.

Spikingformer approfondisce lo stesso punto in una linea di ricerca parallela~\cite{zhou2026spikingformer}. Gli autori analizzano le connessioni residue di Spikformer e mostrano che la variante SEW genera valori a più bit lungo i collegamenti diretti, mentre una Membrane Shortcut (MS) preserva ingressi binari per gli operatori successivi. Combinando questa soluzione con la SSA, il modello mira a conservare insieme comunicazione spike-driven e capacità di modellazione globale.

Il confronto mette in luce un compromesso che accompagna l'intera evoluzione degli Spiking Transformer. Semplificare l'attenzione in modo aggressivo ne riduce il costo, ma impoverisce le relazioni che il modello può rappresentare. Conservare una self-attention più globale offre maggiore espressività, ma obbliga a intervenire altrove per preservare l'efficienza spike-driven. La storia di queste architetture non va letta, allora, come una successione di modelli via via migliori. È piuttosto l'esplorazione di uno spazio definito da almeno tre assi, la capacità relazionale, la scalabilità e l'aderenza al paradigma spike-driven, lungo i quali ciascun lavoro sceglie la propria posizione.

\subsection{La struttura gerarchica}
\label{subsec:struttura_gerarchica}

Resta aperta una questione di altro genere, che riguarda il modo in cui scalare l'architettura stessa. Un Vision Transformer tradizionale mantiene spesso la stessa risoluzione lungo gran parte dell'encoder, e molte delle prime architetture Spiking Transformer seguono la stessa strategia. Spikformer, per esempio, usa uno Spiking Patch Splitting (SPS) per convertire l'ingresso in $N$ token di dimensione $D$, elaborati poi da una sequenza di blocchi che ne conservano il numero e la dimensione (figura~\ref{fig:scala_unica_gerarchia}a).

La soluzione è semplice, ma si discosta dalla struttura gerarchica che caratterizza molte CNN e numerosi Vision Transformer moderni. In una rete gerarchica le prime fasi operano su molte posizioni con pochi canali, mentre quelle profonde lavorano su pochi token descritti da feature più ricche. Questa progressione presenta due vantaggi. Preserva il dettaglio spaziale nelle prime fasi mentre costruisce rappresentazioni via via più astratte, e riduce il numero dei token proprio dove la dimensione delle feature, e con essa il costo dei blocchi Transformer, tende a crescere.

\begin{figure}[t]
    \centering
    \includegraphics[width=\textwidth]{images/cap1/cap1_gerarchia.pdf}
    \caption[Struttura a scala unica e struttura gerarchica]{Struttura a scala unica e struttura gerarchica, secondo le descrizioni originali. I rettangoli rappresentano le feature, con l'altezza proporzionale al lato della griglia dei token e la larghezza al numero dei canali, mentre il colore distingue i blocchi con SSA da quelli con Q-K Token Attention (QKTA). (a)~In Spikformer tutti i blocchi operano sugli stessi $N$ token. (b)~In QKFormer ogni modulo SPEDS riduce la griglia e aumenta i canali. La dimensione temporale $T$ resta invariata lungo entrambe le reti.}
    \label{fig:scala_unica_gerarchia}
\end{figure}

Da questa esigenza nasce QKFormer~\cite{zhou2024qkformer}. Gli autori individuano due limiti congiunti delle architetture precedenti, la complessità quadratica della SSA rispetto al numero dei token e l'elevato costo di memoria che deriva dal combinare la dimensione spaziale con la simulazione temporale. Per superarli introducono una struttura gerarchica a tre stadi, collegati da moduli di Spiking Patch Embedding with Deformed Shortcut (SPEDS), e un nuovo meccanismo di attenzione, la Q-K Attention (figura~\ref{fig:scala_unica_gerarchia}b). Oltre a chiedersi come rendere più economico il prodotto tra Q, K e V, QKFormer si domanda se tutte e tre le componenti siano davvero necessarie per selezionare l'informazione. La Q-K Attention ne usa soltanto due, Query e Key, e riduce l'attenzione a un vettore binario che decide quali token, o quali canali, lasciar passare, come mostra l'ultima riga della figura~\ref{fig:evoluzione_self_attention_spiking}. La sezione~\ref{sec:qkformer_riferimento} ne descrive il funzionamento in dettaglio.

La tabella~\ref{tab:evoluzione_attenzione_spiking} riassume il percorso. QKFormer ne è il punto di partenza naturale per questa tesi, perché riunisce tre proprietà preziose nel regime compatto studiato più avanti, un'attenzione di costo lineare, una rappresentazione gerarchica e una trasmissione residua che attraversa i cambi di risoluzione.

\begin{table}[t]
\centering
\small
\renewcommand{\arraystretch}{1.2}
\caption[Evoluzione dei meccanismi di attenzione spiking]{Evoluzione dei meccanismi di attenzione spiking rilevanti per il lavoro.}
\label{tab:evoluzione_attenzione_spiking}
\begin{adjustbox}{max width=\textwidth}
\begin{tabular}{@{}>{\raggedright\arraybackslash}p{2.6cm} >{\raggedright\arraybackslash}p{3.3cm} >{\raggedright\arraybackslash}p{2.6cm} >{\raggedright\arraybackslash}p{2.1cm} >{\raggedright\arraybackslash}p{2.1cm} c@{}}
\hline
\textbf{Architettura} & \textbf{Problema affrontato} & \textbf{Operatore} & \textbf{Ambito} & \textbf{Costo} & \textbf{Spike-driven} \\
\hline
Transformer / ViT & Relazioni globali tra token & VSA con softmax & Coppie di token & Quadratica in $N$ & No \\
Spikformer & Rendere spiking la self-attention & SSA con Q, K, V spiking & Coppie di token & Quadratica in $N$ & Parziale \\
Spike-driven Transformer & Eliminare le operazioni non spike-driven e ridurre il costo & SDSA, maschere e addizioni sparse & Selezione di canali & Lineare & Sì \\
Spikingformer & Conservare l'attenzione globale con residui spike-driven & SSA con residui MS & Coppie di token & Quadratica in $N$ & Sì \\
QKFormer & Scalabilità e struttura gerarchica & QKTA o QKCA con SPEDS & Selezione di token o canali & Lineare nella dimensione & Nell'attenzione \\
\hline
\end{tabular}
\end{adjustbox}
\vspace{0.4em}
\begin{flushleft}
\footnotesize
\emph{Nota.} «Spike-driven» indica qui la proprietà architetturale dichiarata dai rispettivi lavori. Il costo effettivo dipende comunque da firing rate, operatori periferici, rappresentazione degli input e piattaforma hardware, e verrà discusso nella sezione~\ref{sec:efficienza_neuromorfica_vincoli}.
\end{flushleft}
\end{table}

Altrettanto importante è osservare quale problema QKFormer non si proponga di risolvere. Le operazioni di attenzione descritte finora costruiscono relazioni tra token o canali all'interno delle feature di un singolo timestep. Data una sequenza $X_1,X_2,\dots,X_T$, l'attenzione applicata all'istante $t$ riceve tipicamente la sola $X_t$, e il collegamento con $X_{t-1}$ resta affidato principalmente agli stati dei neuroni che attraversano il tempo. L'espressione \emph{self-attention} suggerisce, in modo intuitivo ma impreciso, un accesso globale all'intera sequenza. La globalità va invece sempre riferita alla dimensione lungo cui l'operatore agisce. Uno Spiking Transformer può disporre di un'attenzione globale sui token e, insieme, di una comunicazione temporale prevalentemente ricorrente e locale.

Resta da capire se la dinamica neuronale descritta nella sezione~\ref{sec:dinamica_neuronale_memoria} basti a governare anche l'interazione tra $X_t$ e $X_{t-k}$, ed è da questa domanda che muove la sezione~\ref{sec:strategie_modellazione_temporale}. Prima conviene però esaminare in dettaglio il punto di arrivo dell'evoluzione descritta, perché QKFormer determina direttamente la baseline su cui si fonda la parte sperimentale della tesi.


% ==============================================================================
%  QKFormer come riferimento architetturale
% ==============================================================================


\section{QKFormer come riferimento architetturale}
\label{sec:qkformer_riferimento}

La sezione~\ref{subsec:struttura_gerarchica} ha presentato QKFormer come risposta ai limiti di scalabilità delle architetture precedenti~\cite{zhou2024qkformer}. Qui se ne riprende ciò che serve al resto del lavoro, perché da questa architettura discende la baseline compatta del capitolo~\ref{chap:baseline_metodologia}. Contano in particolare tre aspetti, la forma esatta del blocco di attenzione, il peso del patch embedding nella gerarchia e il modo in cui il modello tratta il tempo.

\paragraph{Il blocco con Q-K Token Attention}
\label{subsec:qk_attention}

La Q-K Attention ricava dal proprio ingresso $X$ due sole proiezioni spiking,
\begin{equation}
\label{eq:qk_projections_qkformer}
Q = \operatorname{SN}_{Q} \left( \operatorname{BN}(XW_Q) \right), \qquad K = \operatorname{SN}_{K} \left( \operatorname{BN}(XW_K) \right),
\end{equation}
dove BN indica la batch normalization e $Q,K\in\{0,1\}^{T\times N\times D}$, con $T$ timestep, $N$ token e $D$ canali. Nella Q-K Token Attention (QKTA), la variante adottata dalla baseline, le due proiezioni svolgono compiti distinti. La Query decide quali token lasciar passare, la Key ne trasporta il contenuto. Considerando un solo timestep, si sommano per ciascun token gli spike della Query lungo i canali, e il conteggio attraversa un neurone che lo traduce in una decisione binaria,
\begin{equation}
\label{eq:qkta_vettore_attenzione}
A^{\mathrm{tok}} = \operatorname{SN} \left( \sum_{d=1}^{D} Q_{:,d} \right), \qquad A^{\mathrm{tok}} \in \{0, 1\}^{N \times 1}, \qquad X' = A^{\mathrm{tok}} \odot K.
\end{equation}
Un token la cui Query è abbastanza attiva conserva in uscita la propria Key, gli altri vengono azzerati su tutti i canali, come mostra l'ultima riga della figura~\ref{fig:evoluzione_self_attention_spiking}. La variante Q-K Channel Attention applica lo stesso principio all'altro asse, sommando lungo i token per decidere quali canali conservare, ma non viene impiegata in questo lavoro.

Conviene chiarire anche che cosa questa operazione non faccia. Ogni componente di $A^{\mathrm{tok}}$ dipende soltanto dalla Query del token corrispondente, perché le proiezioni $XW_Q$ e $XW_K$ agiscono sui canali di ciascun token separatamente. Diversamente da VSA e SSA, la QKTA non confronta quindi token diversi, ma ne seleziona alcuni in base alla loro attività. Le relazioni tra posizioni diverse restano affidate alle convoluzioni dei patch embedding, su scala locale, e alla SSA dell'ultimo stadio, su scala globale. Il capitolo~\ref{chap:riallocazione_budget} riesamina questa scelta nel regime compatto.

Il guadagno è nel costo. L'attenzione si riduce a un vettore di $N$ elementi, e gli autori riportano una complessità $\mathcal{O}(D)$ in tempo e $\mathcal{O}(N)$ in spazio, contro $\mathcal{O}(N^2D)$ e $\mathcal{O}(N^2+ND)$ della SSA~\cite{zhou2024qkformer}. La somma degli spike ha inoltre una varianza molto più contenuta del prodotto $Q_sK_s^{\top}V_s$, e il fattore di scala $\beta$ può essere eliminato. Questi valori, però, riguardano la sola operazione di attenzione. Le matrici $Q$ e $K$ vanno comunque calcolate e occupano $\mathcal{O}(ND)$, e attorno alla maschera restano proiezioni, normalizzazioni e neuroni. Un'attenzione di costo trascurabile non implica quindi un blocco di costo trascurabile.

Al mascheramento segue infatti una proiezione lineare, che ricombina le feature dei token selezionati,
\begin{equation}
\label{eq:qk_proiezione_lineare_uscita}
X'' = \operatorname{SN} \left( \operatorname{BN} \left( \operatorname{Linear}(X') \right) \right),
\end{equation}
e il blocco completo conserva la struttura di un encoder Transformer,
\begin{align}
\label{eq:blocco_qkformer_qka}
X'_l &= \operatorname{QKA}(X_{l-1}) + X_{l-1}, \\
\label{eq:blocco_qkformer_smlp}
X_l &= \operatorname{SMLP}(X'_l) + X'_l,
\end{align}
dove QKA indica l'intero modulo di attenzione e SMLP uno spiking MLP~\cite{zhou2024qkformer}. Le somme residue avvengono a valle dei neuroni, secondo lo schema SEW della figura~\ref{fig:residui_sew_ms}a. L'attenzione opera quindi con maschere e addizioni, ma ciò che circola tra i blocchi contiene piccoli interi, e gli strati pesati che li ricevono eseguono moltiplicazioni. Per questo i conteggi di questo lavoro distinguono MAC e AC in base all'operando e non al tipo di strato (sezione~\ref{sec:caratterizzazione_computazionale}).

\paragraph{Gerarchia e patch embedding}
\label{subsec:gerarchia_patch}

La struttura a tre stadi è quella della figura~\ref{fig:scala_unica_gerarchia}b. Poiché a ogni passaggio i token diminuiscono e i canali aumentano, un unico tipo di attenzione non è ugualmente adatto a tutta la rete. Gli autori adottano perciò un'integrazione mista, con la QKTA nei primi stadi, dove i token sono molti, e la SSA nell'ultimo, dove sono pochi e il costo quadratico pesa meno. Le prove di ablazione sostengono la scelta. Su CIFAR10-DVS la combinazione di QKTA e SSA raggiunge l'$84{,}0\%$, mentre la sola QKTA si ferma all'$80{,}7\%$~\cite{zhou2024qkformer}.\footnote{Gli autori segnalano nel repository ufficiale un errore nell'ordine di reshape e LIF della proiezione SSA, con possibile sovrastima dei risultati ImageNet. I punteggi dell'articolo vanno quindi distinti da una riproduzione verificata dell'implementazione corretta~\cite{qkformerImplementationNotice}.} La baseline eredita questa sequenza, e il capitolo~\ref{chap:riallocazione_budget} verifica se la SSA resti necessaria anche con una capacità molto inferiore.

Un collegamento residuo ordinario presuppone che ingresso e uscita abbiano le stesse dimensioni, ma il patch embedding che separa due stadi cambia sia il numero dei token sia quello dei canali. Il modulo SPEDS risolve il problema inserendo una proiezione leggera $W_d$ sul ramo residuo,
\begin{equation}
\label{eq:speds_formulazione}
Y = \mathcal{F}(X, \{W_i\}) + \operatorname{SN}(W_d X),
\end{equation}
dove $\mathcal{F}$ è il ramo principale, una sequenza di convoluzioni, normalizzazioni, neuroni e max-pooling, e $W_d$ una convoluzione $1\times 1$ con passo maggiore di uno, che adegua in un solo colpo risoluzione e canali~\cite{zhou2024qkformer}. Da qui il nome di \emph{deformed shortcut}, perché il collegamento diretto non può più essere un'identità e va «deformato» per portare l'ingresso nello spazio dell'uscita. Inserito in Spikformer al posto del modulo originale, SPEDS porta l'accuratezza su CIFAR10-DVS dall'$80{,}90\%$ all'$82{,}20\%$~\cite{zhou2024qkformer}.

La gerarchia ha però un rovescio che conta per questa tesi. Rendendo economica l'attenzione e concentrando la riduzione di scala in moduli dedicati, sposta una parte del budget verso il patch embedding, che lavora sulle mappe a risoluzione più alta. Nella baseline compatta il primo patch embedding ospita oltre l'$80\%$ dello stato persistente, e i due patch embedding insieme circa due terzi dei parametri (sezione~\ref{subsec:bilancio_baseline}). Da questa osservazione parte l'analisi del capitolo~\ref{chap:riallocazione_budget}.

Per i dati neuromorfici gli autori impiegano una configurazione ridotta, con $1{,}50$ milioni di parametri, una larghezza massima di 256 canali e blocchi distribuiti sui tre stadi secondo lo schema $\{0,1,1\}$, nel quale il primo stadio si riduce al solo patch embedding. Con 16 timestep raggiunge il $98{,}6\%$ su DVS128 Gesture e l'$84{,}0\%$ su CIFAR10-DVS~\cite{zhou2024qkformer}. La baseline compatta scende a circa un terzo di questi parametri (tabella~\ref{tab:confronto_qkformer}), e in un regime simile la riduzione non può essere meccanica. Componenti che in QKFormer occupano una frazione modesta del budget possono diventare dominanti, e il peso di ciascuno va rivalutato quando larghezza e profondità diminuiscono.

\paragraph{Il ruolo del tempo}
\label{subsec:qkformer_tempo}

Nelle equazioni precedenti $Q$ e $K$ hanno una dimensione $T$, e si potrebbe pensare che l'attenzione operi congiuntamente sui timestep. Non è così. Gli autori precisano che il tempo resta separato soltanto negli strati di neuroni, mentre negli altri operatori viene fuso con la dimensione del batch~\cite{zhou2024qkformer}. Per questi operatori ogni timestep è un campione a sé, e a ciascun istante $t$ l'attenzione lavora sulla sola rappresentazione di quel passo,
\begin{equation}
\label{eq:attenzione_istante_t}
A^{\mathrm{tok}}_t = \mathcal{A}(Q_t).
\end{equation}
Il passato può influenzare questa decisione soltanto attraverso lo stato dei neuroni, compreso quello che produce il vettore, e non compare alcun termine in cui $Q_t$ venga messa a confronto con $Q_{t-1}$ o con istanti precedenti. QKFormer possiede dunque una struttura fortemente gerarchica nello spazio, ma nessuna gerarchia corrispondente lungo il tempo. Ciò non significa che il modello ignori l'informazione temporale, ma che l'interazione temporale non è l'oggetto del suo meccanismo di attenzione.

Per questa tesi la distinzione è essenziale. QKFormer fornisce un backbone in cui l'elaborazione spaziale e di canale è già stata progettata con cura per essere gerarchica ed efficiente, e consente così di porre una domanda più specifica di quella affrontata dal lavoro originale.
\begin{center}
\emph{Una volta portata l'architettura in un regime fortemente compatto, la sola memoria LIF è il modo più efficace di spendere il budget destinato alla dimensione temporale?}
\end{center}
Il lavoro originale non offre elementi per rispondere, e la risposta va cercata sperimentalmente. La baseline compatta nasce da una scelta conservativa. Mantiene i principi di QKFormer con fedeltà sufficiente a costituire un riferimento riconoscibile, ma ne riduce drasticamente la scala. Soltanto dopo averne caratterizzato il comportamento diventa possibile decidere quali componenti preservare, quali semplificare e quali sostituire. La sezione successiva esamina intanto le strategie proposte in letteratura per rendere esplicito il legame tra istanti diversi.


% ==============================================================================
%  Strategie per la modellazione temporale esplicita
% ==============================================================================


\section{Strategie per la modellazione temporale esplicita}
\label{sec:strategie_modellazione_temporale}

La sezione~\ref{subsec:memoria_relazioni} ha distinto la persistenza di una traccia del passato dall'accesso a ciò che è accaduto in un istante preciso, e la sezione~\ref{subsec:qkformer_tempo} ha mostrato che in QKFormer il passato raggiunge il presente soltanto attraverso lo stato dei neuroni. Il limite non è esclusivo di questa architettura, e Shi et al. hanno cercato di quantificarlo con il Temporal Interaction Coefficient (TIC)~\cite{shi2026multidelay}. Il coefficiente guarda a quanto l'uscita dell'attenzione al passo $q$ sia sensibile agli ingressi dei passi precedenti, e misura con un'entropia quanto questa sensibilità sia distribuita su più istanti. Su CIFAR10-DVS, per SSA, SDSA e QKTA, la dipendenza risulta fortemente concentrata sul passo corrente, mentre i passi lontani contribuiscono pochissimo. Il TIC descrive la forma della dipendenza e non la sua utilità, ma rende misurabile un'osservazione altrimenti qualitativa.

Le due strategie esaminate di seguito affiancano alla memoria dei neuroni un accesso esplicito al passato, e la figura~\ref{fig:confronto_integrazione_passato} le mette a confronto. Entrambe contano il tempo in passi, quindi il significato fisico delle loro distanze dipende dal clock dell'input (sezione~\ref{subsec:discretizzazione_clock}).

\begin{figure}[t]
    \centering
    \includegraphics[width=\textwidth]{images/cap1/cap1_accesso_passato.pdf}
    \caption[Tre modalità di accesso al passato]{Tre modalità di accesso al passato, con il tempo in orizzontale e il passo corrente $t$ a destra. In tutti i casi lo stato dei neuroni collega ogni passo al successivo. (a)~Con la sola memoria neuronale il passato arriva al presente compresso in un unico stato, e il peso di ciascun istante dipende soltanto dalla distanza, come indica l'intensità del colore. (b)~L'attenzione a blocchi confronta direttamente le feature degli istanti che appartengono allo stesso blocco, qui con $B=4$. (c)~Il mixing a ritardi legge alcuni istanti scelti, ciascuno con il proprio peso, qui con i ritardi di 1, 2, 4 e 8 passi del mixer usato in questo lavoro. Lo spessore dei collegamenti è illustrativo.}
    \label{fig:confronto_integrazione_passato}
\end{figure}

\paragraph{Attenzione spazio-temporale a blocchi}
\label{subsec:statten}

STAtten estende la SSA a più timestep~\cite{lee2025spiking}. Anziché calcolare l'attenzione separatamente per ciascun istante, raggruppa $B$ timestep consecutivi e dispone token e istanti lungo un unico asse, come se ogni coppia formata da una posizione e da un istante fosse un token a sé. Con $T=4$ e $B=2$, per esempio, vengono elaborati insieme i timestep $(0,1)$ e $(2,3)$. L'informazione del primo elemento di un blocco non deve più sopravvivere nello stato neuronale fino al secondo, perché le feature di entrambi partecipano alla stessa operazione. Di conseguenza l'operazione non è causale all'interno del blocco, dato che ogni istante vede anche quelli che lo seguono.

Procedere a blocchi ha due motivazioni. Una correlazione sull'intera sequenza richiederebbe di conservare tutte le feature temporali, con un costo di memoria elevato. Inoltre il numero di neuroni attivi dopo il prodotto tra Query, Key e Value diminuisce con la distanza temporale, sicché correlare istanti molto lontani diventa poco produttivo~\cite{lee2025spiking}. La dimensione $B$ definisce così un orizzonte di interazione esplicita, oltre il quale la relazione torna a essere affidata ai neuroni. Sui dati neuromorfici gli autori usano $B=4$ con 16 timestep, e integrata in una configurazione di QKFormer a quattro blocchi STAtten guadagna un punto di accuratezza su CIFAR10-DVS. Il prezzo è che le feature di più timestep devono essere disponibili insieme, un costo che non compare nel conteggio dei parametri e che in un modello compatto pesa in modo particolare.

\paragraph{Mixing mediante ritardi}
\label{subsec:mixing_ritardi}

Il Multi-Delay Mixer (MD-Mixer) segue una strada diversa. Non apre l'attenzione a più timestep, ma modifica le feature che la raggiungono~\cite{shi2026multidelay}. Per ciascun canale $i$ considera $K$ ritardi $d_i^{(1)},\dots,d_i^{(K)}$ e altrettanti coefficienti, tutti appresi durante l'addestramento, e calcola
\begin{equation}
\label{eq:md_mixer_combinazione}
\widetilde X_{t,i} = \sum_{k=1}^{K} \alpha_i^{(k)} X_{t-d_i^{(k)},i}.
\end{equation}
Ogni canale attinge così a più istanti della propria storia, e canali diversi possono imparare ritardi diversi. Il cambiamento rispetto al LIF è qualitativo. Nell'equazione~\eqref{eq:lif_espansione_temporale} è la distanza temporale a fissare il peso di ogni contributo, attraverso il decadimento, mentre qui è la rete a stabilire quanta importanza attribuire a ciascun ritardo, e un contributo lontano raggiunge il presente senza dover sopravvivere a tutti gli aggiornamenti intermedi del potenziale.

Nell'architettura completa il mixer prende il posto della proiezione lineare di Key e Value, mentre la Query resta allineata all'istante corrente. L'attenzione continua a lavorare un passo alla volta, ma su feature che contengono già una parte del passato. Integrato in una configurazione di QKFormer a due blocchi, MD-Mixer guadagna poco più di un punto di accuratezza su CIFAR10-DVS, contro meno di mezzo punto di STAtten nello stesso confronto. Sembra contare soprattutto la possibilità di scegliere quale passato leggere. In un esperimento di ablazione su Spike-driven Transformer, sullo stesso dataset, ritardi casuali e fissi portano circa mezzo punto, mentre ritardi appresi quasi due punti e mezzo~\cite{shi2026multidelay}. Poiché uno spostamento discreto non è derivabile rispetto al ritardo, durante l'addestramento ogni ritardo è una combinazione morbida di più spostamenti, che una temperatura decrescente rende via via più netta fino a una scelta discreta.

È l'approccio più vicino al mixer causale a più ritardi del modello di riferimento di questa tesi (capitolo~\ref{chap:riallocazione_budget}), che ne condivide l'idea di fondo ma se ne discosta in tre punti. Agisce nei patch embedding anziché su Key e Value, combina i canali tra loro con una matrice per ciascun ritardo anziché trattarli separatamente, e usa ritardi fissi di 1, 2, 4 e 8 passi, uguali per tutti i canali.

% TODO: eventuale parte aggiuntiva della 1.5 sugli operatori che confrontano passato e presente
% (correlatori di Hassenstein-Reichardt, termini bilineari). Per ora sono trattati nel Capitolo 5,
% sezione 5.1. Decidere a campagna conclusa; il paragrafo seguente va adattato di conseguenza.

Attenzione a blocchi e mixing a ritardi affiancano dunque la memoria neuronale con esigenze diverse di conservazione delle feature. Per confrontarli occorre guardare, oltre che alle prestazioni, a parametri, operazioni e stato persistente, le grandezze introdotte nella sezione successiva.


% ==============================================================================
%  Efficienza neuromorfica e vincoli hardware
% ==============================================================================

\section{Efficienza neuromorfica e vincoli hardware}
\label{sec:efficienza_neuromorfica_vincoli}

Un modello spiking viene spesso definito efficiente perché comunica mediante spike binari e presenta attivazioni sparse. Queste proprietà sono il presupposto di una computazione a basso consumo, ma non la garantiscono. Il vantaggio si materializza soltanto quando l'assenza di attività si traduce in lavoro effettivamente evitato, e dipende quindi anche dall'hardware su cui il modello viene eseguito. Questo lavoro non realizza i modelli su un dispositivo, ma caratterizza ogni architettura con grandezze che ne anticipano il costo. La sezione le introduce e ne chiarisce i limiti.

\paragraph{Dalle operazioni alla stima energetica}
\label{subsec:operazioni_energia}

La sezione~\ref{subsec:ssa} ha già introdotto la differenza tra una moltiplicazione seguita da un accumulo (MAC), necessaria quando l'attivazione è un valore a più bit, e un semplice accumulo (AC), che basta quando è uno spike binario. Il risparmio dipende però da quanti spike vengono emessi. Il firing rate $\rho$ misura la frazione di neuroni attivi, mediata sui passi. Una rete con $\rho=0{,}05$ attiva in media il $5\%$ dei neuroni a ogni timestep, mentre una con $\rho=0{,}60$ resta formalmente spiking, ma offre molte meno occasioni di saltare le operazioni associate agli zeri.

Il legame tra attività e costo si esprime con il numero di operazioni sinaptiche (\emph{synaptic operations}, SOP). Se uno strato $l$ richiederebbe $F_l$ operazioni in forma densa, in forma spiking ne esegue circa
\begin{equation}
\label{eq:stima_sop}
\operatorname{SOP}_l \approx \rho_l \, T \, F_l,
\end{equation}
perché le operazioni vengono ripetute a ogni timestep, ma soltanto la frazione corrispondente agli spike emessi dà luogo a calcolo. Il profilo di questo lavoro riporta separatamente il conteggio potenziale, in cui ogni possibile spike è presente, e quello pesato per l'attività osservata (sezione~\ref{sec:caratterizzazione_computazionale}).

QKFormer adotta questa impostazione e traduce i conteggi in una stima dell'energia, assegnando a ciascuna operazione un costo di riferimento~\cite{zhou2024qkformer},
\begin{equation}
\label{eq:energia_snn_semplificata}
E_{\mathrm{SNN}} \approx E_{\mathrm{MAC}}N_{\mathrm{MAC}} + E_{\mathrm{AC}}N_{\mathrm{AC}}.
\end{equation}
Nei valori di riferimento una MAC costa circa cinque volte un AC, e il risultato dipende quindi molto da come si attribuiscono le operazioni alle due voci. Nel calcolo di QKFormer le MAC provengono dal solo primo strato, che codifica l'ingresso in spike, e tutto il resto conta come accumulo. Il conteggio di questo lavoro è più severo, perché considera MAC ogni operazione con un operando non binario, comprese quelle che ricevono le somme residue descritte nella sezione~\ref{subsec:qk_attention}. Le stime dei due lavori non vanno perciò confrontate direttamente.

La formula va letta per quello che è, un modello aritmetico e non la misura del consumo di un sistema fisico. Restano fuori, tra l'altro, gli accessi in memoria, il trasferimento dei dati e la logica di controllo, e due modelli con la stessa stima possono comportarsi in modo molto diverso una volta realizzati. In questa tesi la stima serve a confrontare architetture a parità di input, e non viene mai presentata come energia misurata.

\paragraph{Stato persistente e buffer temporali}
\label{subsec:stato_buffer}

I pesi non sono l'unica memoria che un modello spiking occupa. Ogni neurone LIF deve conservare il proprio potenziale di membrana da un passo al successivo, e un'esecuzione in streaming deve leggerlo, aggiornarlo e riscriverlo a ogni timestep. Questo stato persistente segue il numero dei neuroni, e non quello dei pesi. Una convoluzione con pochi parametri, applicata a una mappa ad alta risoluzione, richiede un potenziale per ogni posizione e per ogni canale. Per questa ragione, nella baseline compatta, il primo patch embedding concentra la maggior parte dello stato (sezione~\ref{subsec:gerarchia_patch}) pur contenendo meno di un quinto dei parametri.

Una capacità temporale esplicita può accrescere ulteriormente questo costo. Un mixer che legge fino a $d_{\max}$ passi indietro deve conservare in un buffer le feature di quei passi, e occupa circa $d_{\max}$ volte lo spazio di una singola feature, mentre un'attenzione a blocchi deve avere disponibili insieme le feature di $B$ timestep. Per questo gli elementi di stato e i buffer temporali vengono riportati sempre accanto ai parametri. Un modulo temporale con pochi pesi può rivelarsi poco conveniente, se costringe a conservare molte feature storiche.

Ne discende la regola progettuale che guida il resto del lavoro. La capacità di un modello va giudicata per il luogo in cui viene spesa, non soltanto per la sua quantità. Ridurre i canali delle prime feature map può liberare pochi parametri, ma ridurre in misura molto maggiore i potenziali di membrana e le operazioni ripetute a ogni timestep. Allo stesso modo, un piccolo mixer temporale su feature a bassa risoluzione aggiunge pochi parametri e richiede un buffer ridotto, e può risultare più conveniente di un'attenzione temporale sulle feature ad alta risoluzione. In questo consiste un'allocazione della capacità consapevole dell'hardware (\emph{hardware-aware capacity allocation}), che il capitolo~\ref{chap:riallocazione_budget} applica alla baseline.

\paragraph{Dalla stima all'hardware}
I contatori descrivono il modello, non la macchina che lo esegue. Una GPU non è ottimizzata per la sparsità, perché le ampie regioni inattive dei tensori spiking non vengono sfruttate, e il continuo trasferimento dei dati tra memoria e unità di calcolo aggiunge un costo considerevole~\cite{karamimanesh2025fpga}. Il tempo di inferenza su GPU dice quindi poco sull'efficienza neuromorfica di un modello. L'hardware dedicato cerca invece di rendere visibile la sparsità. L'acceleratore di Li et al. per Spike-driven Transformer opera sugli indirizzi degli spike attivi, con addizioni e confronti anziché moltiplicazioni~\cite{li2024sparseacc}, e quello di Du et al. su FPGA non impiega alcun blocco DSP~\cite{du2024spike}. Il vantaggio dura però finché i circuiti che gestiscono la sparsità costano meno delle operazioni che permettono di evitare~\cite{karamimanesh2025fpga}.

Su una FPGA il costo si esprime infine in risorse del dispositivo, cioè logica (LUT), registri (FF), memoria (BRAM) e blocchi aritmetici (DSP), oltre che in latenza e potenza. Queste grandezze richiedono una sintesi che esula da questo lavoro. I conteggi qui descritti ne sono al più un'anticipazione, e le prospettive di una futura realizzazione vengono riprese nelle conclusioni.

Si chiude così il percorso del capitolo, che ha seguito l'informazione dal segnale prodotto dalla event camera alla sua discretizzazione, dalla memoria neuronale all'attenzione spiking, fino alla modellazione temporale esplicita e al suo costo. L'obiettivo non è costruire il modello con il minor numero di operazioni, né quello più accurato a qualunque costo, ma trovare una soluzione in cui l'informazione temporale aggiuntiva valga le risorse necessarie a rappresentarla. Il capitolo successivo applica questi criteri ai modelli effettivamente implementati, tenendo distinte le scelte relative all'input, i meccanismi architetturali e le condizioni di valutazione.
